\section{Classification}
\label{sec:classification}

We now move from predicting a real number to predicting a \emph{category}. As
explained at the beginning of these notes, the general idea does not change: we
pick a class of functions, measure the quality of a prediction with an appropriate
loss, and minimise the empirical risk. What changes is the nature of $y$, and
therefore the loss. Almost everything we built for regression, and in particular
gradient descent and regularisation, carries over untouched.

\subsection{Binary classification}
\label{subsec:binary}

We start with two categories. The data is as before,
\[
(x_1, y_1), \dots, (x_n, y_n) \in \R^d \times \Y ,
\]
but now the label lives in a set of two elements. It will be very convenient to take
\[
\Y = \{-1, +1\} ,
\]
and we keep this convention throughout the section. Think of $+1$ as ``the email is
spam'' and $-1$ as ``it is not''.

\paragraph{Linear classifiers.}
We keep the linear model $\langle w, x \rangle$, but its output is a real number,
whereas we now need to produce a sign. The natural fix is to threshold it: the
predictor associated with $w \in \R^d$ is
\[
f_w(x) = \sign \big( \langle w, x \rangle \big) \in \{-1, +1\} .
\]
The quantity $\langle w, x \rangle$ is called the \emph{score} of $x$. Geometrically,
the set where the score vanishes,
\[
H_w = \{\, x \in \R^d \ : \ \langle w, x \rangle = 0 \,\} ,
\]
is a hyperplane through the origin (with the intercept trick below, an arbitrary
hyperplane), and $f_w$ answers according to the side of $H_w$ on which $x$ falls. The
hyperplane $H_w$ is called the \emph{decision boundary}, and $w$ is orthogonal to it.

\begin{remark}[The intercept, again]
Exactly as in \Cref{subsec:ridge} and before it, an affine decision boundary
$\langle w, x \rangle + b = 0$ is obtained by appending a constant coordinate $1$ to
each input and an intercept $b$ to the weights. We therefore keep writing $\langle w,
x\rangle$ without loss of generality.
\end{remark}

\paragraph{The margin.}
A key quantity of this section is the product of the label and the score.

\begin{definition}[Margin]
The \emph{margin} of the predictor $w$ on the sample $(x_i, y_i)$ is
\[
m_i(w) \coloneqq y_i \, \langle w, x_i \rangle \ \in \R .
\]
\end{definition}

The margin packages the correctness of a prediction into a single signed number:
\begin{itemize}
    \item $m_i > 0$ means $\langle w, x_i\rangle$ and $y_i$ have the same sign, i.e.\
    the point is \emph{correctly classified};
    \item $m_i < 0$ means the point is \emph{misclassified};
    \item $|m_i|$ measures how far the point sits from the decision boundary, so a
    large positive margin means ``correct, and confidently so''.
\end{itemize}
This is precisely why the labels $\pm 1$ are pleasant: they let a single expression
$y_i \langle w, x_i\rangle$ carry both the prediction and whether it was right. Every
loss we meet below is a function of the margin alone.

\begin{exercise}[$\star$]
Prove that the $| y_i \langle \frac{w}{\Vert w \Vert}, x_i \rangle |$ corresponds to the distance $\min_{x \in H_w} \Vert x - x_i \Vert$ which corresponds to the distance of $x_i$ to $H_w$.
\end{exercise}

\begin{remark}[Labels in $\{0, 1\}$]
Many references, and most software, encode the two classes as $\tilde y \in \{0, 1\}$
rather than $y \in \{-1, +1\}$. Nothing of substance changes: the two conventions are
exchanged by the affine bijection
\[
\tilde y = \frac{1 + y}{2} \in \{0, 1\} ,
\qquad\text{equivalently}\qquad
y = 2 \tilde y - 1 \in \{-1, +1\} .
\]
The $\{0,1\}$ convention is the natural one when the label is read as an
\emph{indicator} of membership in class $1$, which is the form in which most software
expects it; the $\{-1,+1\}$ convention is the natural one when the label is read as a
\emph{sign}, which is what makes the margin, and hence this whole section, work. We
will use $\pm 1$ throughout and translate when useful.
\end{remark}

\subsection{Choosing a loss: from the ideal one to convex surrogates}
\label{subsec:class-losses}

\paragraph{What we would like to minimise.}
The quantity we actually care about is the proportion of mistakes. This corresponds
to the \emph{$0/1$ loss}, which we write directly as a function of the margin $m$:
\[
\ell_{0/1}(m) = \mathbf{1}\{ m \leq 0 \}
= \begin{cases} 1 & \text{if the point is misclassified,} \\ 0 & \text{otherwise,}\end{cases}
\]
whose empirical risk
\[
L_{0/1}(w) = \frac{1}{n} \sum_{i=1}^n \mathbf{1}\big\{ y_i \langle w, x_i \rangle \leq 0 \big\}
\]
is exactly the \emph{misclassification rate} on the training set. Unfortunately it is hopeless to
minimise directly. Since it is piecewise constant, we cannot use gradient methods, its gradient is zero wherever it is defined. Gradient descent, started anywhere, does not move.
Minimising it is in fact NP-hard in general.

\paragraph{Convex surrogates.}
The standard remedy is to replace $\ell_{0/1}$ by a \emph{surrogate loss}: a function
$\ell(m)$ of the margin that is convex and dominates the $0/1$ loss,
possibly after multiplication by a positive constant, so that making it small forces
the misclassification rate to be small too. Here are the classical choices.
\begin{itemize}
    \item \emph{Square loss}, $\ell(m) = (1 - m)^2$. This is nothing but least squares
    applied to the labels $\pm 1$, and it is the one loss of the list with a closed
    form. Its defect is that it is increasing again for $m > 1$: it \emph{penalises
    points that are classified too well}, which is not what we want.
    \item \emph{Hinge loss}, $\ell(m) = \max(0, 1 - m)$. It asks each point to be
    correctly classified with a margin of at least $1$, and is perfectly happy beyond
    that. Combined with an $\ell_2$ penalty it gives the \emph{support vector
    machine}. It is convex but not differentiable at $m = 1$.
    \item \emph{Logistic loss}, $\ell(m) = \ln\big(1 + e^{-m}\big)$. Convex, smooth,
    strictly decreasing, and never exactly zero. This is the one we develop below.
    \item \emph{Exponential loss}, $\ell(m) = e^{-m}$. Convex and smooth, but it
    penalises misclassified points enormously, which makes it sensitive to outliers.
    It is the loss underlying boosting.
\end{itemize}
All of them are decreasing functions of the margin (but the square loss only for $m \leq 1$,
which is precisely its problem), and this is the essential point: they all push each
$m_i$ to be as large as possible, that is, they all reward classifying correctly and
confidently.


\subsection{Logistic regression}
\label{subsec:logistic}

\paragraph{Why this surrogate and not another.}
We focus here on the 
\emph{logistic loss}
\[
\ell(m) = \ln\big( 1 + e^{-m} \big).
\]
This is by far the most used surrogate in practice and it is what essentially every classification
neural network is trained on. It combines all the nice following properties:
\begin{itemize}
    \item It is \emph{convex}, so the empirical risk has no spurious local minima.
    it is the very point of replacing $\ell_{0/1}$.
    \item It is \emph{strictly decreasing}, so a point is never penalised for being
    classified \emph{too} well. 
    \item It grows only \emph{linearly} as $m \to -\infty$, since $\ln(1 + e^{-m}) =
    -m + \ln(1 + e^{m}) \approx -m$. A single badly misclassified point therefore
    contributes a bounded slope and cannot hijack the whole objective. The exponential
    loss fails this badly: $e^{-m}$ blows up, and one outlier can dominate the sum.
\end{itemize}
Note also its behaviour on the other side:
$\ell(m) \to 0$ as $m \to +\infty$ but $\ell(m) > 0$ for every finite $m$. This means that contrary to the square loss, $\ell$ does not have a finite minimiser over $\R$.

\paragraph{The logistic risk.}
Plugging the linear score into the empirical risk with the logistic loss, learning
amounts to minimising over $w \in \R^d$ the function
\[
L(w) = \frac{1}{n} \sum_{i=1}^n \ell\big( m_i(w) \big)
= \frac{1}{n} \sum_{i=1}^n \ln\Big( 1 + \exp\big( - y_i \langle w, x_i \rangle \big) \Big) .
\]
This is the objective of \emph{logistic regression}, and it is the one we minimise for
the rest of this section. Two features of it are already visible from the loss alone.
Being an average of convex functions of $w$, it is convex, so we are again in the
comfortable situation of \Cref{sec:least-squares}. But every term is strictly positive,
so $L(w) > 0$ for every $w$, and the infimum of $L$ may not be attained.

\paragraph{A convenient notation.}
Differentiating $\ell$ produces the same expression over and over, so we give it a
name. Let
\[
\sigma(m) = \frac{1}{1 + e^{-m}} , \qquad m \in \R ,
\]
called the \emph{sigmoid} (or logistic) function. It increases from $0$ to $1$, with
$\sigma(0) = 1/2$, and satisfies the two identities
\[
\sigma(-m) = 1 - \sigma(m) ,
\qquad
\sigma'(m) = \sigma(m)\big(1 - \sigma(m)\big) = \sigma(m)\, \sigma(-m).
\]
It is related to the loss by $\ell(m) = -\ln \sigma(m)$, and:
\[
\ell'(m) = - \sigma(-m) \ \in \ (-1, 0) .
\]
So $\sigma(-m)$ is exactly the \emph{steepness} of the logistic loss at margin $m$:
close to $1$ when $m \ll 0$, close to $0$ when $m \gg 0$.




\begin{remark}[The same objective with $\{0,1\}$ labels]
With the convention $\tilde y \in \{0,1\}$ introduced in \Cref{subsec:binary}, and
writing $\sigma_i = \sigma(\langle w, x_i\rangle)$ for the predicted probability of
class $1$, the formula found in textbooks and libraries is
\[
L(w) = - \frac{1}{n} \sum_{i=1}^n \Big[ \tilde y_i \ln \sigma_i + (1 - \tilde y_i) \ln (1 - \sigma_i) \Big] ,
\]
usually called the \emph{binary cross-entropy}. It is the logistic risk in
disguise. Since $\tilde y_i$ is either $0$ or $1$, the bracket is a switch: it keeps
$\ln \sigma_i = \ln \sigma(\langle w, x_i \rangle)$ when $\tilde y_i = 1$, that is
when $y_i = +1$, and it keeps $\ln(1 - \sigma_i) = \ln \sigma(-\langle w, x_i
\rangle)$ when $\tilde y_i = 0$, that is when $y_i = -1$. In both cases what survives
is $\ln \sigma\big( y_i \langle w, x_i \rangle \big) = \ln \sigma(m_i)$, and the
minus sign in front of the sum turns it into $\ell(m_i) = - \ln \sigma(m_i)$. So this
is exactly the objective above. 
\end{remark}

\begin{exercise}[$\star$]
Prove the last claim: substituting $\tilde y_i = (1 + y_i)/2$ and $\sigma_i =
\sigma(\langle w, x_i\rangle)$ into the cross-entropy formula gives back
$\frac1n \sum_i \ln(1 + e^{-y_i \langle w, x_i\rangle})$.
\end{exercise}

\paragraph{Gradient and Hessian.}
Since we will have to minimise $L$ by an iterative method, we need its gradient. Write
$m_i = m_i(w) = y_i \langle w, x_i \rangle$ for readability.

\begin{prop}[Derivatives of the logistic risk]
\label{prop:logistic-grad}
The function $L$ is differentiable on $\R^d$, with
\[
\nabla L(w) = - \frac{1}{n} \sum_{i=1}^n \sigma(-m_i) \, y_i \, x_i
= - \frac{1}{n} \sum_{i=1}^n \frac{y_i \, x_i}{1 + e^{y_i \langle x_i, w \rangle}} ,
\]
and
\[
\nabla^2 L(w) = \frac{1}{n} \sum_{i=1}^n \sigma(m_i)\,\sigma(-m_i)\, x_i x_i^\top
= \frac{1}{n} X^\top D_w X \ \succeq \ 0 ,
\]
where $D_w = \diag\big( \sigma(m_i) \sigma(-m_i) \big)_{1 \leq i \leq n} \in \R^{n
\times n}$ is diagonal with positive entries.
In particular $L$ is convex. If moreover $\rank(X) = d$, then $\nabla^2 L(w) \succ 0$
for every $w$ and $L$ is strictly convex.
\end{prop}

\begin{proof}
The scalar function $\ell(m) = \ln(1 + e^{-m})$ has derivative
\[
\ell'(m) = \frac{-e^{-m}}{1 + e^{-m}} = \frac{-1}{1 + e^{m}} = -\sigma(-m) ,
\]
and the chain rule applied to $w \mapsto m_i(w) = y_i \langle w, x_i\rangle$, whose
gradient is $y_i x_i$, gives the announced $\nabla L$. Differentiating once more,
\[
\nabla_w \big[ \sigma(-m_i) \big] = \sigma'(-m_i) \cdot (-y_i x_i)
= - \sigma(-m_i)\sigma(m_i) \, y_i x_i ,
\]
so that $\nabla^2 L(w) = \frac1n \sum_i \sigma(m_i)\sigma(-m_i) \, y_i^2 \, x_i
x_i^\top$, and $y_i^2 = 1$ because $y_i \in \{-1,+1\}$. Each coefficient
$\sigma(m_i)\sigma(-m_i)$ is $> 0$, so for any $u \in \R^d$,
\[
u^\top \nabla^2 L(w) u = \frac1n \sum_{i=1}^n \sigma(m_i)\sigma(-m_i) \, \langle u, x_i \rangle^2 \geq 0 ,
\]
which is the convexity. If $\rank(X) = d$ then $\ker(X) = \{0\}$, so $u \neq 0$ forces
$\langle u, x_i \rangle \neq 0$ for at least one $i$, and the sum is $> 0$.
\end{proof}

\begin{remark}[Reading the gradient]
The formula for $\nabla L$ is worth staring at. A gradient \emph{descent} step moves
$w$ along $-\nabla L(w)$, that is, along
\[
\frac{1}{n} \sum_{i=1}^n \underbrace{\sigma(-m_i)}_{\text{weight} \, \in \, (0,1)} \ \underbrace{y_i x_i}_{\text{direction}} .
\]
Each sample pushes $w$ in the direction $y_i x_i$, which is exactly the direction that
increases its own margin $m_i = y_i \langle w, x_i \rangle$: every point asks to be
classified better. What differs between points is the weight they are given, and that
weight is the steepness $\sigma(-m_i) = -\ell'(m_i)$ met above. Badly classified points
($m_i \ll 0$) carry a weight close to $1$ and dominate the step; comfortably classified
points ($m_i \gg 0$) carry a weight close to $0$ and are almost ignored. Logistic
regression therefore spends nearly all its effort on the points near the decision
boundary or on the wrong side of it, and essentially none on the easy ones.
\end{remark}

\paragraph{No closed form.}
Here is the essential practical difference with least squares. There, setting $\nabla
L(w) = 0$ gave the \emph{linear} system $X^\top X w = X^\top y$, which we could solve
in closed form. Here, $\nabla L(w) = 0$ reads
\[
\sum_{i=1}^n \frac{y_i \, x_i}{1 + e^{y_i \langle w, x_i \rangle}} = 0 ,
\]
in which the unknown $w$ appears inside an exponential. This is a system of $d$
transcendental equations, and it has \emph{no closed-form solution}: there is no
analogue of the normal equations, and no formula such as $(X^\top X)^{-1} X^\top y$ to
write down. We are therefore in the second of the two situations listed at the
beginning of \Cref{subsec:gd}, and we do the only thing we can: we minimise $L$
\emph{iteratively}. Concretely,
\[
w_{t+1} = w_t - \gamma\, \nabla L(w_t)
\qquad\text{or}\qquad
w_{t+1} = w_t - \gamma\, \nabla \ell_{i_t}(w_t) ,
\]
that is, exactly the gradient descent and stochastic gradient descent of
\Cref{subsec:gd}, with the gradient of \Cref{prop:logistic-grad} plugged in. Since $L$
is convex and smooth, gradient descent with a small enough step size will minimise the training loss.

\subsection{Separable and non-separable data}
\label{subsec:separable}

\begin{definition}[Linearly separable data]
The training set is said to be \emph{linearly separable} if there exists $w \in \R^d$ such that
\[
y_i \langle w, x_i \rangle > 0 \qquad \text{for every } i = 1, \dots, n ,
\]
that is, if some hyperplane classifies every training point correctly. Otherwise it is
\emph{non-separable}.
\end{definition}

Whether or not the data is separable completely changes the behaviour of $L$.

\paragraph{Separable data: the minimum is at infinity.}

\begin{prop}[No minimiser in the separable case]
\label{prop:sep-no-min}
Suppose the data is linearly separable. Then
\[
\inf_{w \in \R^d} L(w) = 0 ,
\]
but the infimum is \emph{not attained}: $L(w) > 0$ for every $w \in \R^d$. Moreover
every minimising sequence $(w_k)$ (i.e. such that $L(w_k) \to 0$) satisfies $\Vert w_k \Vert \to \infty$.
\end{prop}

\begin{proof}
Each term $\ln(1 + e^{-m})$ is strictly positive for every finite $m$, so $L(w) > 0$
for every $w$, and the infimum cannot be attained at any finite point.

Let $u$ be a separating direction and set $\rho \coloneqq \min_i y_i \langle u, x_i
\rangle > 0$, the smallest margin it achieves. Walking along the ray
$t u$ with $t > 0$, every margin is multiplied by $t$:
\[
L(t u) = \frac1n \sum_{i=1}^n \ln\big( 1 + e^{-t\, y_i \langle u, x_i\rangle} \big)
\ \leq \ \ln\big( 1 + e^{-t \rho} \big) \ \xrightarrow[t \to +\infty]{} \ 0 ,
\]
so $\inf L = 0$. Finally, if some minimising sequence had a bounded subsequence, it
would have a convergent one, $w_{k_j} \to \bar w$, and continuity of $L$ would give
$L(\bar w) = 0$, contradicting $L > 0$. Hence $\Vert w_k \Vert \to \infty$.
\end{proof}

The picture to keep in mind is this: once a separating direction $u$ is found, scaling
it up multiplies all the margins by $t$ without ever changing a single prediction,
since $\sign(\langle t u, x\rangle) = \sign(\langle u, x \rangle)$. The classifier is
already perfect, but the loss keeps rewarding more confidence, and it can always buy
more by making $\Vert w \Vert$ larger. So the optimisation problem has no solution:
\emph{the minimum sits at infinity}, and gradient descent, faithfully, sends $\Vert
w_t \Vert$ to infinity. The training loss decreases to $0$ while the training accuracy has
been $100\%$ for a long time.

\begin{remark}[Which direction does it diverge in?]
If $\Vert w_t \Vert \to \infty$, the only thing that still matters is the
\emph{direction} $w_t / \Vert w_t \Vert$, since that is what determines the
predictions. And the direction does converge. This is the exact analogue, for the
logistic loss, of the implicit bias of gradient descent we proved in
\Cref{subsec:gd} for least squares: the loss alone does not pick a separating
hyperplane, but the algorithm does. We record the result without proof.
\end{remark}

\begin{fact}[Implicit bias of GD on separable data, admitted]
On linearly separable data, gradient descent on the logistic loss with a small enough
constant step size satisfies $\Vert w_t \Vert \to \infty$,
while
\[
\frac{w_t}{\Vert w_t \Vert} \ \xrightarrow[t \to \infty]{} \ \frac{\hat u}{\Vert \hat u \Vert} ,
\qquad
\hat u \in \argmax_{\Vert u \Vert \leq 1} \ \min_{1 \leq i \leq n} \ y_i \langle u, x_i \rangle .
\]
That is, the direction converges to the \emph{maximum-margin} separator, the one
staying as far as possible from the closest training points. This is the solution of
the hard-margin support vector machine, obtained here for free, without ever writing
that problem down.
\end{fact}

\paragraph{Non-separable data: a minimum exists.}
When the two classes overlap, no $w$ can classify everything correctly, the loss can
no longer be driven to $0$. 

\begin{prop}[Existence of a minimiser]
\label{prop:nonsep-min}
Suppose that for every $w \neq 0$ there exists a datapoint such that  $y_i \langle w, x_i
\rangle < 0$, i.e.\ that every candidate direction \emph{strictly} misclassifies at
least one point. Then $L$ is coercive and attains its minimum.
\end{prop}

\begin{proof}
The function $h(u) \coloneqq \min_i y_i \langle u, x_i \rangle$ is continuous (a
minimum of finitely many linear functions) and, by assumption, $h(u) < 0$ on the unit
sphere, which is compact. Hence $\rho \coloneqq - \max_{\Vert u \Vert = 1} h(u) > 0$:
every unit direction misclassifies some point with margin at most $-\rho$.

Now take any $w \neq 0$ and write $w = t u$ with $t = \Vert w \Vert$ and $\Vert u
\Vert = 1$. Picking an index $i$ with $y_i \langle u, x_i \rangle \leq -\rho$ and
keeping only that term,
\[
L(w) \ \geq \ \frac1n \ln\big( 1 + e^{t \rho} \big) \ \geq \ \frac{t \rho}{n}
\ \xrightarrow[\ \Vert w \Vert \to \infty\ ]{} \ \infty ,
\]
so $L$ is coercive. A continuous coercive function attains its minimum. 
\end{proof}


\paragraph{Weight regularisation.}
Just as ridge regression cured the non-uniqueness of least squares in the
over-parametrised regime, adding an $\ell_2$ penalty cures the ``minimum at infinity''
problem here. Consider
\[
L_\lambda(w) = \frac{1}{n} \sum_{i=1}^n \ln\big( 1 + e^{- y_i \langle w, x_i \rangle} \big)
\ + \ \frac{\lambda}{2} \Vert w \Vert^2 , \qquad \lambda > 0 .
\]
Being a sum of a convex function and a strictly convex one, $L_\lambda$ is strictly
convex; and $L_\lambda(w) \geq \frac{\lambda}{2}\Vert w \Vert^2 \to \infty$, so it is
coercive. By the argument used for ridge in \Cref{subsec:ridge}, it therefore admits a
\emph{unique} minimiser, whatever the data, separable or not, and whatever the
relative sizes of $n$ and $d$. The penalty is precisely what forbids the escape to
infinity: buying more confidence now costs $\frac{\lambda}{2}\Vert w \Vert^2$, and at
some finite point the purchase is no longer worth it. As before, there is still no
closed form, and one minimises $L_\lambda$ by gradient descent, whose gradient is
$\nabla L(w) + \lambda w$.


\subsection{Multiclass classification}
\label{subsec:multiclass}

We finally treat $K \geq 2$ categories, so that
\[
y_i \in \{1, 2, \dots, K\} ,
\]
for instance $K = 10$ when recognising which digit is written in an image.

\paragraph{One score per class.}
A single score is no longer enough: we compute \emph{one score per class}, using one
weight vector per class. The parameter is now a collection
\[
W = (w_1, \dots, w_K), \qquad w_k \in \R^d ,
\]
which we may equally see as a matrix $W \in \R^{d \times K}$, and the scores of an
input $x$ are the $K$ numbers $\langle w_1, x \rangle, \dots, \langle w_K, x \rangle$.
The prediction is the class with the largest score,
\[
f_W(x) = \argmax_{1 \leq k \leq K} \ \langle w_k, x \rangle .
\]

\paragraph{The multiclass margin.}
The reasoning of \Cref{subsec:class-losses} repeats verbatim, and again it is the
$0/1$ loss that we would like to minimise and cannot. Writing $s_k = \langle w_k, x_i
\rangle$ for the scores of the sample $i$, the prediction is correct exactly when the
true class wins, that is, when $s_{y_i} > s_k$ for every $k \neq y_i$, equivalently
when
\[
\max_{k \neq y_i} \ \big( s_k - s_{y_i} \big) \ < \ 0 .
\]
This quantity plays the role that $-m_i$ played in the binary case: it is negative
exactly when the point is correctly classified, and the $0/1$ loss is the indicator
that it is non-negative. It is once more non-convex and flat, so once more we look for
a convex surrogate.

\paragraph{Softening the maximum.}
The obstruction is the maximum, which is convex but not smooth. The standard fix is to
replace it by the \emph{log-sum-exp} function, which is a smooth approximation of the
maximum: for any reals $a_1, \dots, a_K$,
\[
\max_{k} a_k \ \leq \ \ln \sum_{k=1}^K e^{a_k} \ \leq \ \max_{k} a_k + \ln K .
\]
The left inequality holds because the sum dominates its largest term, the right one
because the sum is at most $K$ times its largest term. So log-sum-exp is the maximum,
up to an additive constant of at most $\ln K$, and it is convex and $C^\infty$.
Applying it to $a_k = s_k - s_{y_i}$ (which now includes the index $k = y_i$,
contributing $e^0 = 1$) gives the loss
\[
\ell_i(W) = \ln \sum_{k=1}^K \exp\big( s_k - s_{y_i} \big)
= \ln \Big( 1 + \sum_{k \neq y_i} \exp\big( s_k - s_{y_i} \big) \Big) ,
\]
and the empirical risk, after factoring $e^{-s_{y_i}}$ out of the sum:
\[
\boxed{\
L(W) = \frac{1}{n} \sum_{i=1}^n \left[ - \langle w_{y_i}, x_i \rangle
+ \ln \sum_{k=1}^K \exp\big( \langle w_k, x_i \rangle \big) \right] .
\ }
\]
This is the \emph{cross-entropy} (or \emph{softmax}) loss. The first term pushes up the
score of the correct class, the second pushes down the scores of all classes; their
balance is what makes the correct class win. The middle expression makes the kinship
with the binary case plain: it is $\ln(1 + \sum_{k \neq y} e^{s_k - s_y})$, against
$\ln(1 + e^{-m})$ for two classes. This is by some distance the most used loss in
modern machine learning: it is the objective at the end of essentially every
classification network. Evaluating it on a computer requires a little care, since the
exponentials overflow easily; see \Cref{subsec:numerics}.

\paragraph{The softmax weights.}
Differentiating log-sum-exp produces normalised exponentials, which we name once and
for all:
\[
p_k(x) \coloneqq \frac{\exp\big( \langle w_k, x \rangle \big)}{\sum_{k'=1}^K \exp\big( \langle w_{k'}, x \rangle \big)} ,
\qquad k = 1, \dots, K ,
\]
called the \emph{softmax weights}. They are positive, they sum to $1$, and the largest
score gets the largest weight. They are the multiclass analogue of the steepness
$\sigma(-m)$ of the binary case, and indeed for $K = 2$ they reduce to it.

\begin{prop}[Gradient of the cross-entropy]
\label{prop:softmax-grad}
For each $k \in \{1, \dots, K\}$,
\[
\nabla_{w_k} L(W) = \frac{1}{n} \sum_{i=1}^n \Big( p_k(x_i) - \mathbf{1}\{ y_i = k \} \Big) \, x_i .
\]
Moreover $L$ is convex in $W$.
\end{prop}

The factor $p_k(x_i) - \mathbf{1}\{y_i = k\}$ compares the weight the
softmax currently puts on class $k$ with the weight we would like it to put there,
namely $1$ for the true class and $0$ for all the others. If the sample is already
classified with a comfortable margin, the softmax weights are already close to that
target, the factor is nearly zero and the sample contributes almost nothing to the
step; if it is misclassified, the factor is close to $\pm 1$ and the sample pushes
hard. This is the exact multiclass counterpart of the reading we gave of the binary
gradient. Exactly as in the binary case, too, setting all the $\nabla_{w_k} L$ to zero
gives a system of transcendental equations with no closed-form solution, and we
minimise by gradient descent.

\begin{exercise}[$\star\star$]
Prove \Cref{prop:softmax-grad}. Start by showing that the gradient of the second term,
$\nabla_{w_k} \ln \sum_{k'} e^{\langle w_{k'}, x\rangle}$, is equal to $p_k(x)\, x$,
and then differentiate the first term, which is nonzero only when $k = y_i$.
\end{exercise}

\begin{remark}[The parametrisation is redundant]
Adding the same vector $v \in \R^d$ to \emph{every} $w_k$ leaves the softmax
unchanged, since $e^{\langle w_k + v, x\rangle} = e^{\langle v, x\rangle} e^{\langle
w_k, x\rangle}$ and the common factor cancels in the ratio. So $L(W)$ is constant
along all such directions, and a minimiser is \emph{never} unique: as in the
over-parametrised least-squares problem of \Cref{subsec:overparam}, the set of
minimisers contains a whole subspace of directions the loss cannot see. One fixes this
either by regularising (adding $\frac{\lambda}{2}\sum_k \Vert w_k \Vert^2$, which
selects the representative of smallest norm) or by fixing $w_K = 0$ by convention.
\end{remark}

\paragraph{Recovering binary classification.}
Take $K = 2$, with class $1$ standing for the label $y = +1$ and class $2$ for $y =
-1$. The middle form of the loss makes the computation immediate: with only one rival
class, the inner sum has a single term, and the loss of a sample of class $1$ is
\[
\ell_i = \ln\Big( 1 + e^{\,s_2 - s_1} \Big)
= \ln\Big( 1 + e^{- \langle w_1 - w_2, \, x_i \rangle} \Big) ,
\]
while for a sample of class $2$ the roles of the indices are swapped and
\[
\ell_i = \ln\Big( 1 + e^{\,s_1 - s_2} \Big)
= \ln\Big( 1 + e^{+ \langle w_1 - w_2, \, x_i \rangle} \Big) .
\]
The two weight vectors therefore enter only through their \emph{difference}, which is
exactly the redundancy of the previous remark: the loss is constant along the
$d$-dimensional subspace $\{(v, v) : v \in \R^d\}$ of the parameter space $\R^{2d}$,
and what survives is the $d$-dimensional quotient, coordinatised by $w_1 - w_2$.
Setting
\[
w \coloneqq w_1 - w_2 \in \R^d ,
\]
both cases become the single expression $\ln(1 + e^{-y_i \langle w, x_i\rangle})$, so
the multiclass cross-entropy with $K = 2$ \emph{is} the logistic risk of
\Cref{subsec:logistic}, written in the variable $w = w_1 - w_2$. Binary logistic
regression is thus not a separate method but the special case $K = 2$ of softmax
regression, and the $\pm 1$ margin formulation is what you obtain after removing the
redundant direction.

\begin{exercise}[$\star$]
Check the analogous statement for the predictions rather than the loss: verify that
$\argmax_{k \in \{1,2\}} \langle w_k, x\rangle = 1$ if and only if $\langle w, x
\rangle > 0$, so that the multiclass rule reduces to $\sign(\langle w, x\rangle)$.
\end{exercise}